\section{Descripción de la Arquitectura General}

El sistema se organiza bajo una arquitectura en capas (\textit{layered
architecture}) que separa las responsabilidades de captura, servicios,
procesamiento, persistencia y presentación. Esta organización responde a la
naturaleza asimétrica de la carga de trabajo: la recepción de solicitudes es
ligera y debe permanecer siempre disponible, mientras que la inferencia de los
modelos de aprendizaje profundo es intensiva en cómputo y su duración es
variable. Separar ambas responsabilidades permite que el servicio de recepción
responda de inmediato y delegue la inferencia a un proceso trabajador
independiente, evitando que una ráfaga de vehículos bloquee la atención de las
solicitudes.

La Figura \ref{fig:arquitectura_general} presenta la arquitectura general del
sistema, en la que cada capa se apoya exclusivamente en las interfaces públicas
de la capa adyacente. El flujo de información es descendente en el sentido de la
percepción ---desde la imagen capturada hacia el registro persistido--- y
ascendente en el sentido de la supervisión, a través de las consultas y de las
notificaciones en tiempo real que el sistema envía al panel de monitoreo.

\begin{figure}[htbp]
    \centering
    \begin{tikzpicture}

        % *** Capas arquitectónicas ***
        \spLayer{capaCaptura}{0,0}{13cm}{1.4cm}{cielo}{Capa de Captura}
        \spLayer{capaServicios}{0,-2.3}{13cm}{1.7cm}{purpura}{Capa de Servicios}
        \spLayer{capaProceso}{0,-4.8}{13cm}{1.7cm}{matblue}
            {Capa de Procesamiento Asíncrono}
        \spLayer{capaDatos}{0,-7.0}{13cm}{1.4cm}{purpura}{Capa de Persistencia}
        \spLayer{capaPresent}{0,-9.2}{13cm}{1.7cm}{rosa}{Capa de Presentación}

        % *** Componentes: captura ***
        \spBlock{camara}{0,0}{4.6cm}{cielo}
            {Cámara de acceso}{Captura del vehículo}

        % *** Componentes: servicios ***
        \spBlock{notif}{-2.0,-2.3}{3.8cm}{purpura}
            {Notificador}{WebSocket}
        \spBlock{api}{2.0,-2.3}{3.8cm}{purpura}
            {API REST}{FastAPI}

        % *** Componentes: procesamiento ***
        \spBlock{cola}{-2.0,-4.8}{3.8cm}{matblue}
            {Cola de tareas}{Redis}
        \spBlock{worker}{2.0,-4.8}{3.8cm}{matblue}
            {Worker de inferencia}{YOLO26 + detector de placa + CRNN}

        % *** Componentes: persistencia ***
        \spBlock{bd}{0,-7.0}{5.4cm}{purpura}
            {PostgreSQL}{Vehículos y sesiones}

        % *** Componentes: presentación ***
        \spBlock{web}{-2.0,-9.2}{3.8cm}{rosa}
            {Aplicación Web}{React}
        \spBlock{operador}{2.0,-9.2}{3.8cm}{rosa}
            {Operador}{Supervisión}

        % *** Flujo principal entre capas ***
        \draw[spflow] (capaCaptura.south) --
            node[right, font=\tiny\itshape, text=black!70] {imagen}
            (capaServicios.north);
        \draw[spflow] (capaServicios.south) --
            node[right, font=\tiny\itshape, text=black!70] {tarea}
            (capaProceso.north);
        \draw[spflow] (capaProceso.south) --
            node[right, font=\tiny\itshape, text=black!70] {resultados}
            (capaDatos.north);
        \draw[spflow] (capaDatos.south) --
            node[right, font=\tiny\itshape, text=black!70] {consultas}
            (capaPresent.north);

        % *** Relaciones internas ***
        \draw[spflow] (cola.east) -- (worker.west);
        \draw[spboth] (web.east) -- (operador.west);

        % *** Notificación en tiempo real ***
        \draw[spflow] (notif.west) -- (-6.9,-2.3) -- (-6.9,-9.2) --
            (web.west);
        \node[above, font=\tiny\itshape, text=black!70]
            at (-5.4,-2.3) {notificación};

    \end{tikzpicture}
    \caption{Arquitectura general en capas del sistema de gestión inteligente de
    parqueaderos.}
    \label{fig:arquitectura_general}
\end{figure}

\subsection{Capa de captura}

Constituye el punto de entrada de información al sistema y está conformada por
las cámaras instaladas en los accesos del parqueadero. Su responsabilidad es
adquirir la imagen del vehículo que aproxima a la barrera y remitirla al
servicio de gestión. El diseño de esta capa se mantiene deliberadamente simple
para favorecer la reutilización de la infraestructura de captura existente y
para desacoplar el sistema de un modelo específico de cámara: basta con que el
dispositivo de captura pueda entregar una imagen en un formato conocido.

\subsection{Capa de servicios}

Expone la interfaz pública del sistema y constituye el único punto de acceso a
sus funcionalidades. La API REST recibe las solicitudes de registro de ingreso y
de salida, valida su estructura y delega el trabajo pesado a la capa de
procesamiento, respondiendo de inmediato con el identificador de la sesión
creada. El notificador mantiene un canal persistente con los clientes web y
propaga las actualizaciones de estado sin que estos deban consultar
periódicamente al servidor. Esta capa no conoce los detalles internos de los
modelos de visión, lo que permite evolucionar el motor de inferencia sin alterar
el contrato expuesto a los clientes.

\subsection{Capa de procesamiento asíncrono}

Es el núcleo analítico del proyecto. La cola de tareas desacopla la recepción de
la solicitud de su ejecución: cada solicitud de registro encola una tarea y el
proceso trabajador la consume cuando dispone de capacidad. El trabajador de
inferencia ejecuta la detección de vehículos, el seguimiento, el reconocimiento
de la matrícula y la clasificación de atributos, y deposita el resultado en la
capa de persistencia. Al residir en un proceso separado, esta capa puede
desplegarse sobre hardware con aceleración gráfica ---incluso en equipos
próximos al punto de captura--- y aprovechar modelos optimizados para reducir su
latencia, sin condicionar la disponibilidad del servicio de recepción.

\subsection{Capa de persistencia}

Almacena el estado del sistema: los vehículos registrados, las sesiones de
estacionamiento con sus marcas temporales de ingreso y salida, y la bitácora de
detecciones producida por el proceso de inferencia. La base de datos relacional
actúa como fuente única de verdad y es manipulada exclusivamente por los
servicios de la capa superior, garantizando que las reglas de integridad se
apliquen en un único punto.

\subsection{Capa de presentación}

Ofrece la interfaz mediante la cual el personal del parqueadero supervisa la
operación. La aplicación web consume la API REST para las consultas históricas y
mantiene una conexión de notificación para recibir las actualizaciones en
tiempo real. Al ser una capa puramente de consumo, puede desplegarse de forma
independiente y ser reemplazada o ampliada sin impacto sobre el resto de la
arquitectura.

\subsection{Flujo de comunicación}

La comunicación entre capas sigue un contrato estricto. La imagen capturada se
remite al servicio de gestión, que crea el registro de la sesión y encola la
tarea de inferencia; el trabajador consume la tarea, ejecuta los modelos y
actualiza la sesión con los resultados obtenidos. Una vez consolidado el
resultado, el notificador propaga el cambio a la aplicación web. Esta secuencia
garantiza que ningún componente acceda directamente a los datos de otro: la base
de datos es manipulada por los servicios, y los modelos de visión no conocen la
estructura del almacenamiento. Dicha separación de responsabilidades simplifica
el mantenimiento, facilita las pruebas por componente y sostiene la escalabilidad
del sistema ante la incorporación de nuevos puntos de acceso.
